\section{Finite admissibility and contextual distinctions}\label{sec:context}

% ledger: C1--C6 (2026-09-13 supplement); these are written deductions, not new Lean claims.
The results above give a concrete instance of a distinction in the Six Birds (SBT)
calculus \cite{Tsi26F4}: a rule can survive composition while its realizations require
successively finer contextual distinctions. This complementary reading proceeds from
a local obstruction to the intrinsic hierarchy, then asks what the construction
suggests about descriptions of collective systems. The deductions below are written
proofs; no new Lean verification is claimed.

\emph{A distinction that cannot be erased.}
Let $C(x,y,z)=(y,z,x)$. Since $C\carrier=\carrier$, the nine pairs
$(C^a\solid,C^b\solid+(1,1,1))$, for $a,b\in\{0,1,2\}$, have the same ordered
pair of bare carriers. Their feature profiles match precisely when $a=b$.
The exact checker \file{paper/scripts/contextual_frames.py} reconstructs the $24$
opposed feature centres in each pair from the solid's rational feature records:
every off-diagonal pair mismatches all $24$ coefficients and has an interior overlap.
For example, $\solid$ and $C\solid+(1,1,1)$ both contain
$(5/4,11/8,1)$ in their interiors: there the root's role $188$ protrudes in $+z$
by $12/10000$, and the neighbour's role $4$ protrudes in $-z$ by $5/10000$.
The diagonal pairs are rotated copies of the $000$/central subdivision pair.
In SBT terminology this is a \emph{descent obstruction}: admissibility is not a
function of the erased carrier description. Testing the neighbour against each
of the three root frames requires retaining all three neighbour-frame distinctions.

\emph{Two different closure questions.}
SBT's \emph{sufficiency closure} (F7 of \cite{Tsi26F4}) asks whether each declared
readout factors through an observation. It does not assert that an operation
preserves legal configurations. Here the latter is the separate geometric statement
$D(X_Q)\subseteq X_Q$, where $X_Q$ is the space of tilings by $\solid$
(\cref{thm:coarsening}). This permits repeated recognition and restandardization
under the same rule. The observations of position in that hierarchy behave differently.

\begin{proposition}[phase distinctions]\label{prop:phases}
Fix a registered tiling $T$ in integer coordinates. Its level-$n$ parent anchors
lie in a unique coset $\alpha_n(T)+2^n\Z^3$, with
$\alpha_{n+1}(T)\equiv\alpha_n(T)\pmod{2^n}$. The observations
\[
 q_n^T(v)=v-\alpha_n(T)\pmod{2^n},\qquad v\in\Z^3,
\]
have $8^n$ values. Writing $\ker q=\{(v,w):q(v)=q(w)\}$,
\[
 \ker q_n^T=\{(v,w):w-v\in2^n\Z^3\},\qquad
 \ker q_{n+1}^T\subsetneq\ker q_n^T,\qquad
 \bigcap_{n\ge0}\ker q_n^T=\Delta,
\]
where $\Delta=\{(v,v):v\in\Z^3\}$.
\end{proposition}

\begin{proof}
The common-parent parity argument of \cref{thm:coarsening}, applied at every
level and expressed in the original coordinates, gives the cosets. Each parent
has child $000$ at its own anchor, so the anchor sets are nested and the cosets
are coherent. Reduction modulo $2^n$ is onto and has the displayed kernel.
The pair $v,v+2^ne_1$ lies in that kernel but not the next; their intersection
is $\Delta$ because $\bigcap_n2^n\Z^3=\{0\}$.
\end{proof}

The phases are read from the intrinsic hierarchy, locally at each fixed level
with a neighbourhood growing with the level. In particular, for $p\in\Z^3$,
$\alpha_n(T+p)=\alpha_n(T)+p\pmod{2^n}$. A period must therefore vanish modulo
every $2^n$, recovering \cref{thm:noperiod}. This uses the global partitions,
not exhaustion by one tile's ancestors. Equal phases need not give equal patches:
$8^n$ counts phase states, not patch types or entropy. Each next level requires
eight times as many phase states, while each fixed phase observation has the
lawful translation update $q_n^T(v+u)=q_n^T(v)+u$.

\emph{Exact contextual sufficiency.}
Let $\ell:\Z^3\to A$ be the $168$-label cell encoding of $T$
(\cref{sec:limitperiodic}). It is faithful: a cell's position, native frame
and native cell identify its owner's whole pose. Declare the readouts to be
all translated labels, $v\mapsto\ell(v+u)$ for $u\in\Z^3$. Their common
equivalence, the \emph{predictive quotient} for this spatial probe, satisfies
\[
 v\equiv w\ \Longleftrightarrow\
 \forall u\in\Z^3,\ \ell(v+u)=\ell(w+u)
 \ \Longleftrightarrow\ w-v\in\Per(\ell).
\]
Indeed, put $u=z-v$; the reverse implication is the definition of a period.
Faithfulness gives $\Per(\ell)=\Per(T)=\{0\}$. Thus any observation sufficient
for all these readouts must distinguish every probe origin. This consequence
holds for any faithful aperiodic lattice labeling; Chair44 additionally supplies
the intrinsic phase hierarchy organizing such distinctions.

This sufficiency requirement concerns the declared probes. It excludes neither
finite descriptions nor useful coarse models. The coherent phases give a $2$-adic
coordinate, but need not determine the whole tiling; the hull and diffraction
questions retain the scope stated in \cref{sec:limitperiodic}.

\emph{Collective descriptions and their limits.}
These results suggest a broader question: which relationships must a description
retain for a collective to support its declared operations? An inventory of parts
can identify arrangements that behave differently when placed in a further context.
The appropriate criterion is therefore relative to a specified family of contexts
$\mathcal C$ and their observable outcomes:
\[
 q(x)=q(y)\quad\Longrightarrow\quad
 \operatorname{Outcome}(C[x])=\operatorname{Outcome}(C[y])
 \quad\text{for every }C\in\mathcal C.
\]
Here the outcome includes whether the composition is admissible at all. The
coarsest equivalence satisfying this criterion identifies exactly those states
that all declared contexts fail to distinguish. It specifies the relational
information required by those tests, rather than a complete microscopic inventory.

In SBT terms, a separating context witnesses a descent obstruction. Enlarging the
context family can only refine its observational equivalence, but strict refinement
at arbitrarily large depths requires a separate argument. Recursive composition
alone does not supply one. For Chair44, \cref{prop:phases} supplies that argument:
each phase level identifies positions separated by some nonzero translations,
while no such identification survives every level. Independently, coarsening
preserves the geometric admissibility law that makes the hierarchy repeatable.

The philosophical consequence is precise and limited: a finite, recursively stable
law can require a realization to preserve relational distinctions at arbitrarily
large scales. The infinite arrangement carries those distinctions; the law
constrains their organization. This suggests a question for other collective
systems: which permitted interaction reveals information erased by a proposed
description, and does retaining that information suffice after further composition?
An answer requires identifying that system's own states, operations and separating
contexts. Chair44 provides a geometric instance in which both lawful recursion and
unbounded contextual differentiation can be established.
